
\documentclass[12pt, a4paper]{article}

% --- Pakiety ---
\usepackage[utf8]{inputenc}
\usepackage[T1]{fontenc}
\usepackage[polish]{babel}
\usepackage{amsmath, amssymb}
\usepackage{graphicx}
\usepackage{float}
\usepackage{geometry}
\usepackage{hyperref}
\usepackage{booktabs}
\usepackage{caption}
\usepackage{listings}
\usepackage{xcolor}
\usepackage{parskip}

\geometry{
  a4paper,
  left=2.5cm,
  right=2.5cm,
  top=2.5cm,
  bottom=2.5cm
}

% --- Zakaz dzielenia wyrazów ---
\hyphenpenalty=10000
\exhyphenpenalty=10000
\emergencystretch=3em

\hypersetup{
  colorlinks=true,
  linkcolor=blue,
  urlcolor=blue,
  citecolor=blue
}

% --- Zmiana nazwy float figure: Rysunek -> Wykres ---
\addto\captionspolish{\renewcommand{\figurename}{Wykres}}

% --- Styl listingów ---
\lstset{
  language=Python,
  basicstyle=\ttfamily\small,
  keywordstyle=\color{blue}\bfseries,
  commentstyle=\color{gray}\itshape,
  stringstyle=\color{orange!80!black},
  numbers=left,
  numberstyle=\tiny\color{gray},
  numbersep=8pt,
  frame=single,
  breaklines=true,
  breakatwhitespace=false,
  captionpos=b,
  showstringspaces=false,
  inputencoding=utf8,
  extendedchars=true,
  literate={ą}{{\k{a}}}1 {ę}{{\k{e}}}1 {ó}{{\'{o}}}1
            {ś}{{\'{s}}}1 {ł}{{\l{}}}1 {ź}{{\'{z}}}1
            {ż}{{\.z}}1 {ć}{{\'{c}}}1 {ń}{{\'{n}}}1
            {Ą}{{\k{A}}}1 {Ę}{{\k{E}}}1 {Ó}{{\'{O}}}1
            {Ś}{{\'{S}}}1 {Ł}{{L}}1,
}

% --- Strona tytułowa ---
\title{
  \textbf{Metody Obliczeniowe w Nauce i Technice} \\[0.5em]
  \large Laboratorium 4 \textemdash{} Efekt Rungego i interpolacja
}
\author{Mateusz Stoch}
\date{7 kwietnia 2026}

\begin{document}

\maketitle

% ============================================================
\section{Zadanie 1 \textemdash{} Średnia geometryczna odległości węzłów}
% ============================================================

% --- 1.1 Opis zadania ---
\subsection{Opis zadania}

Dla trzech rodzajów rozkładu węzłów interpolacji na przedziale $[-1,1]$:
\begin{itemize}
  \item węzłów Czebyszewa ($n = 5,\, 10,\, 20$),
  \item węzłów Legendre'a ($n = 5,\, 10,\, 20$),
  \item węzłów równomiernych ($n = 5,\, 10,\, 20$),
\end{itemize}
należy zwizualizować, jak rozkładają się węzły na przedziale, wykreślając
dla każdego węzła $x_j$ jego \emph{średnią geometryczną odległości} do
pozostałych węzłów:
\[
  G_j \;=\; \Bigl(\prod_{k \neq j} |x_j - x_k|\Bigr)^{\!1/(n-1)}.
\]

% --- 1.2 Argumentacja rozwiązania ---
\subsection{Argumentacja rozwiązania}

Średnia geometryczna odległości jest miarą tego, jak daleko od siebie
są węzły w~pobliżu danego punktu. Dla węzłów równomiernych wartość ta
jest niemal stała w~środku przedziału, ale drastycznie spada na
krańcach, ponieważ węzły są tam zagęszczone dokładnie tak samo jak w~środku,
więc odległości do sąsiadów są podobne, ale do odległych węzłów są duże.
Węzły Czebyszewa i~Legendre'a zagęszczają się na brzegach przedziału, co
niweluje efekt Rungego.
\newpage
\begin{lstlisting}[caption={Obliczanie i wizualizacja średniej geometrycznej odległości}]
def plot_geometric_mean_distances(x_nodes, title_suffix='', ax=None):
    n = len(x_nodes)
    geo_means = np.zeros(n)
    for i in range(n):
        dists = np.abs(x_nodes[i] - np.delete(x_nodes, i))
        geo_means[i] = np.prod(dists) ** (1.0 / (n - 1))
    if ax is None:
        fig, ax = plt.subplots()
    ax.plot(x_nodes, geo_means, 'o-', markersize=4)
    ax.set_xlabel('x')
    ax.set_ylabel('Sr. geom. odleglosci')
    ax.set_title(title_suffix)
    return geo_means

def chebyshev_nodes(n, a=-1, b=1):
    j = np.arange(1, n + 1)
    return 0.5*(a+b) + 0.5*(b-a)*np.cos((2*j-1)/(2*n)*np.pi)

def legendre_nodes(n):
    coeffs = np.zeros(n + 1)
    coeffs[n] = 1.0
    return legroots(coeffs)
\end{lstlisting}

% --- 1.3 Wykresy i wyniki ---
\subsection{Wykresy i wyniki}

\begin{figure}[H]
  \centering
  \includegraphics[width=0.85\textwidth]{wykres1.png}
  \caption{Średnia geometryczna odległości węzłów do pozostałych dla
           trzech rodzajów węzłów i~trzech wartości~$n$.}
  \label{fig:geo_means}
\end{figure}

% --- 1.4 Wnioski ---
\subsection{Wnioski}

\begin{enumerate}
  \item Węzły równomierne dają prawie stałą średnią geometryczną w~środku
        przedziału, ale wartości na jego krańcach są wyraźnie niższe.
        Oznacza to, że kraje przedziału są słabo ``pokryte``, przy
        interpolacji wielomianem Lagrange'a stąd wynika efekt Rungego:
        duże oscylacje na brzegach.

  \item Węzły Czebyszewa i~Legendre'a mają podobny wzorzec, są
        zagęszczone przy krańcach przedziału, co sprawia że średnia
        geometryczna odległości jest tam wyraźnie większa niż w~środku,
        ale za to znacznie bardziej równomiernie rozłożona. To właśnie
        ta własność chroni przed efektem Rungego.

  \item Węzły Czebyszewa i~Legendre'a dają niemal identyczne wyniki,
        obie te rodziny węzłów są splecione z~teorią wielomianów
        ortogonalnych i~obie dają zbliżone gwarancje numeryczne.
\end{enumerate}

% ============================================================
\section{Zadanie 2 \textemdash{} Interpolacja i efekt Rungego}
% ============================================================

% --- 2.1 Opis zadania ---
\subsection{Opis zadania}

Celem zadania jest interpolacja dwóch funkcji:
\[
  f_1(x) = \frac{1}{1 + 25x^2}, \quad x \in [-1, 1]
\]
\[
  f_2(x) = e^{\cos(x)}, \quad x \in [0, 2\pi]
\]
za pomocą trzech metod:
\begin{itemize}
  \item wielomianu Lagrange'a z~węzłami równomiernie rozmieszczonymi,
  \item kubicznych funkcji sklejanych z~węzłami równomiernie rozmieszczonymi,
  \item wielomianu Lagrange'a z~węzłami Czebyszewa.
\end{itemize}

Część (a) dotyczy wykresu porównawczego dla $f_1$ z~$n = 12$ węzłami.
Część (b) bada jak zmienia się norma błędu przy zwiększaniu $n$ od~4
do~50.

% --- 2.2 Argumentacja rozwiązania ---
\subsection{Argumentacja rozwiązania}

\subsubsection{Interpolacja barycentryczna Lagrange'a}

Zamiast klasycznego wzoru Lagrange'a (który jest niestabilny), implementuje
się interpolację barycentryczną. Polega ona na obliczeniu wag $w_j$:
\[
  w_j = \frac{1}{\displaystyle\prod_{k \neq j}(x_j - x_k)},
\]
a następnie ewaluacji:
\[
  p(x) = \frac{\displaystyle\sum_j \frac{w_j}{x - x_j}\, y_j}
              {\displaystyle\sum_j \frac{w_j}{x - x_j}}.
\]
Ta forma jest numerycznie znacznie bardziej stabilna.
\newpage
\begin{lstlisting}[caption={Interpolacja barycentryczna Lagrange'a}]
def lagrange_barycentric(x_nodes, y_nodes, x_eval):
    n = len(x_nodes)
    w = np.ones(n)
    for j in range(n):
        for k in range(n):
            if k != j:
                w[j] /= (x_nodes[j] - x_nodes[k])
    x_eval = np.asarray(x_eval, dtype=float)
    result = np.zeros_like(x_eval)
    for i, xi in enumerate(x_eval):
        diffs = xi - x_nodes
        zero_idx = np.where(diffs == 0)[0]
        if len(zero_idx) > 0:
            result[i] = y_nodes[zero_idx[0]]
        else:
            terms = w / diffs
            result[i] = np.dot(terms, y_nodes) / np.sum(terms)
    return result
\end{lstlisting}

\subsubsection{Węzły Czebyszewa}

Węzły Czebyszewa na przedziale $[a, b]$ są zdefiniowane jako:
\[
  x_j = \frac{a+b}{2} + \frac{b-a}{2}\cos\!\left(\frac{2j-1}{2n}\pi\right),
  \quad j = 1, \ldots, n.
\]
Minimalizują one maksymalny błąd interpolacji i~eliminują efekt Rungego.

\subsubsection{Funkcje sklejane}

Do interpolacji użyto \texttt{CubicSpline} z~biblioteki SciPy, który
dopasowuje kawałkowe wielomiany stopnia~3 z~zachowaniem ciągłości
i~gładkości na węzłach.

\begin{lstlisting}[caption={Interpolacja kubiczną funkcją sklejaną}]
cs_unif = CubicSpline(x_unif, y_unif)
y_spline_unif = cs_unif(x_dense_unif)
\end{lstlisting}

% --- 2.3 Wykresy i wyniki ---
\subsection{Wykresy i wyniki liczbowe}

\subsubsection{(a) Interpolacja $f_1$ przy $n=12$ węzłach}

\begin{figure}[H]
  \centering
  \includegraphics[width=0.9\textwidth]{wykres2.png}
  \caption{Porównanie metod interpolacji funkcji Rungego $f_1(x)$ przy
           $n = 12$ węzłach.}
  \label{fig:runge_n12}
\end{figure}

Widać wyraźnie, że Lagrange z~węzłami równomiernie rozmieszczonymi mocno
oscyluje przy krańcach przedziału (efekt Rungego). Funkcja sklejana
i~Lagrange z~węzłami Czebyszewa radzą sobie znacznie lepiej.

\subsubsection{(b) Norma błędu vs. liczba węzłów}

Błąd interpolacji mierzony jest jako norma $\|\mathbf{e}\|_2$ na 500
losowo wybranych punktach testowych.

\begin{figure}[H]
  \centering
  \includegraphics[width=0.9\textwidth]{wykres3.png}
  \caption{Norma błędu interpolacji $f_1(x) = 1/(1+25x^2)$ w~zależności
           od liczby węzłów $n$ (skala logarytmiczna na osi~$y$).}
  \label{fig:error_f1}
\end{figure}

\begin{figure}[H]
  \centering
  \includegraphics[width=0.9\textwidth]{wykres4.png}
  \caption{Norma błędu interpolacji $f_2(x) = e^{\cos(x)}$ w~zależności
           od liczby węzłów $n$ (skala logarytmiczna na osi~$y$).}
  \label{fig:error_f2}
\end{figure}

% --- 2.4 Wnioski ---
\subsection{Wnioski}

\begin{enumerate}
  \item Dla funkcji Rungego $f_1(x)$ Lagrange z~węzłami równomiernie
        rozmieszczonymi jest zdecydowanie najgorszy, błąd rośnie wraz
        z~$n$, co jest doskonałym przykładem efektu Rungego. Dodawanie
        kolejnych węzłów zamiast poprawiać interpolacji, tylko ją
        pogarsza.

  \item Lagrange z~węzłami Czebyszewa eliminuje efekt Rungego, błąd
        maleje monotonicznie wraz ze wzrostem $n$. Zmiana wyboru węzłów
        (przy tej samej metodzie interpolacji) całkowicie zmienia
        zachowanie metody.

  \item Kubiczna funkcja sklejana z~węzłami równomiernie rozmieszczonymi
        zachowuje się dobrze dla obu funkcji, błąd maleje monotonicznie
        i~osiąga niskie wartości. Funkcje sklejane działają lokalnie,
        więc efekt Rungego nie może się propagować na cały przedział.

  \item \textbf{Podsumowanie:} Najdokładniejszą metodą (w~sensie
        najszybszego zaniku błędu) okazuje się Lagrange z~węzłami
        Czebyszewa. Funkcja sklejana jest bardzo dobrą i~bezpieczną
        alternatywą, działa dobrze nawet przy równomiernie
        rozmieszczonych węzłach. Lagrange z~węzłami równomiernie
        rozmieszczonymi jest metodą najmniej dokładną dla funkcji
        z~ostrymi zmianami.
\end{enumerate}

% ============================================================
\section{Podsumowanie}
% ============================================================

\begin{enumerate}
  \item \textbf{Efekt Rungego} jest realnym problemem w~interpolacji
        wielomianowej. Przy węzłach równomiernie rozmieszczonych zwiększanie
        liczby węzłów może pogorszyć interpolacji, a~nie ją poprawić.

  \item \textbf{Wybór węzłów ma kluczowe znaczenie.} Węzły Czebyszewa
        (i~podobnie Legendre'a) minimalizują błąd interpolacji i~eliminują
        efekt Rungego. Różnica pomiędzy nimi a~węzłami równomiernie
        rozmieszczonymi widoczna jest wyraźnie w~rozkładzie średniej
        geometrycznej odległości.

  \item \textbf{Funkcje sklejane są bezpiecznym i~praktycznym wyborem.}
        Nawet przy równomiernie rozmieszczonych węzłach dają dobre
        wyniki,działają lokalnie, a zaburzenie w~jednym miejscu nie
        wpływa na cały przedział.

  \item \textbf{Interpolacja barycentryczna Lagrange'a} jest numerycznie
        stabilną implementacją wzoru Lagrange'a i~powinna być preferowana
        nad klasyczną postacią.
\end{enumerate}

\end{document}
